\section{Details on Fairness Notions}
\label{sec:notions-extra}

Although we do not define any new fairness notions in this paper,
we extend some of them to more general settings than they were originally defined for.
For some notions, this extension is not obvious and is based on careful deliberation.
Here we show how we arrived at these extensions and why they make sense.

\subsection{EFX}
\label{sec:notions:efx}

Defining EFX in the fully general setting (non-additive valuations, mixed manna) is tricky.
So, we start with the definition of EFX for additive goods,
and gradually build up to the general definition of EFX from there.
(TODO: Some of this is part of the EEFX journal paper. How to cite that?)

There are actually two competing defintions of EFX for additive goods.
According to the original definition by \cite{caragiannis2019unreasonable},
an allocation $A$ is EFX-fair to agent $i$ if for every other agent $j$,
removing any positively-valued good from $j$'s bundle eliminates $i$'s envy. Formally,
\[ \frac{v_i(A_i)}{w_i} \ge \max_{g \in A_j: v_i(g) > 0} \frac{v_i(A_j \setminus \{g\})}{w_j}. \]

A different definition, often called \EFXZero{},
doesn't require the good $j$ to have a positive value to $i$ \cite{plaut2020almost}.
An allocation $A$ is \EFXZero-fair to agent $i$ if for every other agent $j$,
\[ \frac{v_i(A_i)}{w_i} \ge \max_{g \in A_j} \frac{v_i(A_j \setminus \{g\})}{w_j}. \]

\EFXZero{} is known to be incompatible with PO
based on the following simple example (TODO: cite),
whereas it is not known whether EFX is compatible with PO.

\begin{example}
\label{ex:efx-po}
Consider the fair division instance $([2], \{g_1, g_2, g_3\}, (v_i)_{i=1}^2, (1/2, 1/2))$,
where $v_1$ and $v_2$ are additive, and
\begin{align*}
   v_1(g_1) &= 1, & v_1(g_2) &= 0, & v_1(g_3) &= 10,
\\ v_2(g_1) &= 0, & v_2(g_2) &= 1, & v_2(g_3) &= 10.
\end{align*}
\end{example}

For \cref{ex:efx-po}, for any Pareto-optimal allocation $A$, we have $g_1 \in A_1$ and $g_2 \in A_2$
(otherwise we can transfer $g_1$ to agent 1 or $g_2$ to agent 2 to obtain a Pareto-dominator of $A$).
Then the agent who did not get $g_3$ in $A$ is not \EFXZero-satisfied, although she is EFX-satisfied.

\begin{comment}
For two agents having additive valuations, when the items are goods,
every MMS-fair allocation is EFX (Theorem 4.6 in \cite{caragiannis2019unreasonable}).
(For more than two agents, PMMS implies EFX \cite{caragiannis2019unreasonable}
and MMS implies EEFX \cite{caragiannis2023new}.)
However, this is not true for \EFXZero.
For \cref{ex:efx-po}, each agent's maximin share is 1, so allocation $(\{1\}, \{2, 3\})$
is MMS, but we have already seen that it is not \EFXZero-fair.
\end{comment}

While \EFXZero{} is trivial to extend to non-additive valuations, EFX is not.
This is because every good in $j$'s bundle can have zero value to agent $i$.
(Additionally, $j$'s marginal value over $A_j$ and $A_i$ may also be zero.)
Before we try to address this issue, let's instead jump to the setting of additive mixed manna.

One way to define EFX for mixed manna is this:
agent $i$ is EFX-satisfied by allocation $A$ if for every other agent $j$,
either agent $i$ doesn't envy $j$, or $i$'s envy towards $j$ vanishes after
either removing the least valuable positively-valued item from $j$
or after removing the most valuable negatively-valued item from $i$.
However, we argue that this is not sufficient.

\begin{example}
\label{ex:efx-mixed-manna}
Consier a fair division instance $\Ical$ having 2 agents with equal entitlements,
identical additive valuations, two goods of values $10$ each,
and two chores of values $-9$ each.
Consider an allocation $A$ where all the 4 items are allocated to agent 2.
Then agent 1 would be EFX-satisfied by $A$, even though allocation $B$,
where each agent gets one good and one chore, is fairer.
\end{example}

For goods, EFX is considered one of the strongest notions of fairness,
so we would like it to be a very strong notion for mixed manna too.
The key observation in \cref{ex:efx-mixed-manna} is that in allocation $A$,
we can transfer a set of items (containing one good and one chore) from agent 2 to agent 1
and get a fairer allocation.
This suggests that instead of (hypothetically) removing a single good from $j$
or a single chore from $i$ in the definition of EFX,
we should remove a positively-valued subset of $A_j$
or a negatively-valued subset of $A_i$.

Hence, for additive mixed manna, an allocation $A$ is EFX-fair to agent $i$
if for every other agent $j$, $i$ doesn't envy $j$, or both of the following hold:
\begin{enumerate}
\item $\displaystyle \frac{v_i(A_i)}{w_i} \ge \frac{\max(\{v_i(A_j \setminus S) \mid S \subseteq A_j
    \textrm{ and } v_i(S) > 0\})}{w_j}$.
\item $\displaystyle \frac{\min(\{v_i(A_i \setminus S) \mid S \subseteq A_i
    \textrm{ and } v_i(S) < 0 \})}{w_i} \ge \frac{v_i(A_j)}{w_j}$.
\end{enumerate}

This definition also hints towards how to handle goods with non-additive valuations.
Even if every good in $j$'s bundle has (marginal) value zero to agent $i$,
some subset of $j$'s bundle must have positive (marginal) value.
We replace $v(S) > 0$ by $v_i(S \mid A_i) > 0$ to ensure that
transferring $S$ from $j$ to $i$ leads to an increase in $i$'s valuation,
otherwise we lose compatibility with PO, as the following example demonstrates.

\begin{example}
\label{ex:efx-two-colors}
Consider a fair division instance with 2 agents having equal entitlements and identical valuations.
There are 2 red goods and 4 blue goods. The value of any bundle is $\max(k_r, k_b)$,
where $k_r$ and $k_b$ are the numbers of red and blue goods in the bundle, respectively.
(The valuation function is the rank function of a partition matroid,
and hence, is submodular. (TODO: prove or cite or remove))
\end{example}

For \cref{ex:efx-two-colors}, in any allocation, some agent gets at most 2 blue goods,
and that agent's value for her own bundle can be at most 2.
Also, the maximum value that any bundle can have is 4.
Hence, an allocation is PO iff one agent gets all the blue goods
and the other agent gets all the red goods.
If we want EFX and PO to be compatible for this instance,
we must define EFX such that the allocation obtained by giving
all blue goods to one agent and all red goods to the other agent must be EFX.

We now identify special cases where our definition of EFX (\cref{defn:efx})
is equivalent to well-known definitions of EFX under special cases.

\begin{lemma}
\label{thm:efx-equiv-positive}
In the fair division instance $\Ical \defeq (N, M, (v_i)_{i \in N}, w)$,
if all marginals are positive for agent $i$
(i.e., $v_i(g \mid R) > 0$ for all $R \subseteq M$ and $g \in M \setminus R$),
then $A$ is EFX-fair to agent $i$ iff for all $j \in N \setminus \{i\}$, we have
\[ \frac{v_i(A_i)}{w_i} \ge \max_{g \in A_j} \frac{v_i(A_j \setminus \{g\})}{w_j}. \]
\end{lemma}
\begin{proof}[Proof sketch]
In \cref{defn:efx}, $v_i(A_j \setminus S)$ is maximized by picking
the minimal $S$ such that $v_i(S \mid A_i) > 0$.
Since marginals are positive, $v_i(S \mid A_i) > 0$ for any $|S| \ge 1$.
Hence, in \cref{defn:efx}, we assume $|S| = 1$ \wLoG{}.
\end{proof}

\begin{lemma}
\label{thm:efx-equiv-negative}
In the fair division instance $\Ical \defeq (N, M, (v_i)_{i \in N}, w)$,
if all marginals are negative for agent $i$
(i.e., $v_i(c \mid R) < 0$ for all $R \subseteq M$ and $c \in M \setminus R$),
then $A$ is EFX-fair to agent $i$ iff for all $j \in N \setminus \{i\}$, we have
\[ \max_{c \in A_i} \frac{|v_i(A_i \setminus \{c\})|}{w_i} \le \frac{|v_i(A_j)|}{w_j}. \]
\end{lemma}
\begin{proof}[Proof sketch]
In \cref{defn:efx}, $|v_i(A_i \setminus S)|$ is maximized by picking
the minimal $S$ such that $v_i(S \mid A_i \setminus S) < 0$.
Since marginals are negative, $v_i(S \mid A_i \setminus S) < 0$ for any $|S| \ge 1$.
Hence, in \cref{defn:efx}, we assume $|S| = 1$ \wLoG{}.
\end{proof}

\begin{lemma}
\label{thm:efx-equiv-submod-goods}
In the fair division instance $\Ical \defeq (N, M, (v_i)_{i \in N}, w)$,
if $v_i$ is submodular and all marginals are non-negative for agent $i$
(i.e., $v_i(g \mid R) \ge 0$ for all $R \subseteq M$ and $g \in M \setminus R$),
then $A$ is EFX-fair to agent $i$ iff for all $j \in N \setminus \{i\}$, we have
\[ \frac{v_i(A_i)}{w_i} \ge \frac{\max(\{v_i(A_j \setminus \{g\})
        \mid g \in A_j \textrm{ and } v_i(g \mid A_i) > 0\})}{w_j}. \]
\end{lemma}
\begin{proof}[Proof sketch]
In \cref{defn:efx}, $v_i(A_j \setminus S)$ is maximized by picking
the minimal $S$ such that $v_i(S \mid A_i) > 0$.
We will show that $v_i(S \mid A_i) > 0$ implies $\exists g \in S$
such that $v_i(g \mid A_i) > 0$, so we can assume \wLoG{} that $|S| = 1$.

Let $S \subseteq A_j$ such that $v_i(S \mid A_i) > 0$.
Let $S \defeq \{g_1, \ldots, g_k\}$. Then by submodularity of $v_i$, we get
\[ 0 < v_i(S \mid A_i) = \sum_{t=1}^k v_i(g_t \mid A_i \cup \{g_1, \ldots, g_{t-1}\})
    \le \sum_{t=1}^k v_i(g_t \mid A_i). \]
Hence, $v_i(g \mid A_i) > 0$ for some $g \in S$.
\end{proof}

\begin{lemma}
\label{thm:efx-equiv-submod-chores}
In the fair division instance $\Ical \defeq (N, M, (v_i)_{i \in N}, w)$,
if $v_i$ is submodular and all marginals are non-positive for agent $i$
(i.e., $v_i(c \mid R) \le 0$ for all $R \subseteq M$ and $c \in M \setminus R$),
then $A$ is EFX-fair to agent $i$ iff for all $j \in N \setminus \{i\}$, we have
\[ \frac{\max(\{|v_i(A_i \setminus \{c\})| \mid c \in A_i
    \textrm{ and } |v_i(c \mid A_i \setminus \{c\})| > 0\})}{w_i} \le \frac{|v_i(A_j)|}{w_j}. \]
\end{lemma}
\begin{proof}[Proof sketch]
In \cref{defn:efx}, $|v_i(A_i \setminus S)|$ is maximized by picking
the minimal $S$ such that $v_i(S \mid A_i \setminus S) < 0$.
We will show that $v_i(S \mid A_i \setminus S) < 0$ implies $\exists c \in S$
such that $v_i(c \mid A_i \setminus \{c\}) < 0$, so we can assume \wLoG{} that $|S| = 1$.

Let $S \subseteq A_i$ such that $v_i(S \mid A_i \setminus S) < 0$.
Let $S \defeq \{c_1, \ldots, c_k\}$.
For any $t \in \{0\} \cup [k]$, let $S_t \defeq \{c_1, \ldots, c_t\}$. Then
\begin{align*}
v_i(S \mid A_i \setminus S) &= v_i(A_i \setminus S_0) - v_i(A_i \setminus S_k)
\\ &= \sum_{t=1}^k (v_i(A_i \setminus S_{t-1}) - v_i(A_i \setminus S_t))
\\ &= \sum_{t=1}^k v_i(c_t \mid A_i \setminus S_t).
\end{align*}
By submodularity of $v_i$, we get
\[ 0 > v_i(S \mid A_i \setminus S) = \sum_{t=1}^k v_i(c_t \mid A_i \setminus S_t)
    \ge \sum_{t=1}^k v_i(c_t \mid A_i \setminus \{c_t\}). \]
Hence, for some $c \in S$, we have $v_i(c \mid A_i \setminus \{c\}) < 0$.
\end{proof}

\subsection{MMS}
\label{sec:notions:mms}

\begin{lemma}
\label{thm:pess-is-mms}
For any fair division instance with equal entitlements,
the pessShare of any agent is the same as her maximin share.
\end{lemma}
\begin{proof}
Let $\Ical \defeq ([n], [m], (v_i)_{i=1}^n, w)$ be a fair division instance,
where $w_i = 1/n$ for all $i \in [n]$.
Any agent $i$'s $1$-out-of-$n$ share is the same as $\MMS_i$, so $\pessShare_i \ge \MMS_i$.

Let $\ell/d \le n$ and let $X \in \Pi_d(M)$ such that $v_i(X_j) \le v_i(X_{j+1})$ for all $j \in [d-1]$.
Now let $Y_1$ be the union of the first $\ell$ bundles of $X$,
let $Y_2$ be the union of the next $\ell$ bundles of $X$, and so on.
Add any remaining bundles of $X$ to $Y_n$.
Formally, $Y \in \Pi_n(M)$ where $Y_k \defeq \bigcup_{j=(k-1)\ell+1}^{k\ell} X_j$ for $k \in [n-1]$
and $Y_n \defeq M \setminus \bigcup_{j=1}^{n-1} Y_j$.
Then $Y_n$ contains at least $\ell$ bundles of $X$, since $d \ge \ell n$.
Hence, $v_i(Y_1) \le v_i(Y_2) \le \ldots \le v_i(Y_n)$,
and $v_i(Y_1)$ is agent $i$'s $\ell$-out-of-$d$ share.

By definition of MMS, $v_i(Y_1) \le \MMS_i$.
Hence, for any $\ell$ and $d$ such that $\ell/d \le n$,
agent $i$'s $\ell$-out-of-$d$ share is at most her MMS.
Hence, $\pessShare \le \MMS_i$.
\end{proof}

\subsection{APS}
\label{sec:notions:aps}

\begin{lemma}
\label{thm:aps-optimal-price}
Let $\Ical \defeq ([n], [m], (v_i)_{i=1}^n, w)$ be a fair division instance.
For agent $i$, let $G$ be the set of goods and $C$ be the set of chores,
i.e., $G \defeq \{g \in [m]: v_i(g \mid R) \ge 0 \forall R \subseteq [m] \setminus \{g\}\}$
and $C \defeq \{c \in [m]: v_i(c \mid R) \le 0 \forall R \subseteq [m] \setminus \{g\}\}$
Then for some optimal price vector $\phat \in \mathbb{R}^m$, we have
$\phat_g \ge 0$ for all $g \in G$ and $\phat_c \le 0$ for all $c \in C$.
\end{lemma}
\begin{proof}
Let $p^* \in \mathbb{R}^m$ be an optimal price vector.
Let $\Ghat \defeq \{g \in G: p^*_g < 0\}$ and $\Chat \defeq \{c \in C: p^*_c > 0\}$.
%
The high-level idea is that if we change the price of $\Ghat \cup \Chat$ to 0,
we get potentially better prices.
%
Define $\phat \in \mathbb{R}^m$ as
\[ \phat_j \defeq \begin{cases}
0 & \textrm{ if } j \in \Ghat \cup \Chat
\\ p^*_j & \textrm{ otherwise}
\end{cases}, \]
and let
\[ \Shat \in \argmax_{S \subseteq [m]: \phat(S) \le w_i\phat([m])} v_i(S) .\]
Since $\phat(\Shat \cup \Ghat \setminus \Chat) = \phat(\Shat)$
and $v_i(\Shat \cup \Ghat \setminus \Chat) \ge v_i(\Shat)$,
we can assume without loss of generality that
$\Ghat \subseteq \Shat$ and $\Chat \cap \Shat = \emptyset$.
\[ p^*(\Shat) - w_ip^*([m])
= (\phat(\Shat) - w_i\phat([m])) - (1-w_i)(-p^*(\Ghat)) - w_ip^*(\Chat) \le 0. \]
Hence,
\[ \max_{S \subseteq [m]: p^*(S) \le w_ip^*([m])} v_i(S)
\ge v_i(\Shat) = \max_{S \subseteq [m]: \phat(S) \le w_i\phat([m])} v_i(S), \]
so $\phat$ is also an optimal price vector.
\end{proof}

When all items are goods, by \cref{thm:aps-optimal-price}, we can assume \wLoG{} that
$p$ is non-negative and $p([m]) = 1$. Hence,
\[ \APS_i = \min_{p \in \Delta_m}\;\max_{S \subseteq [m]: p(S) \le w_i} v_i(S). \]

When all items are chores, by \cref{thm:aps-optimal-price}, we can assume \wLoG{} that
$p$ is non-positive and $p([m]) = -1$. Hence,
\[ -\APS_i = \max_{q \in \Delta_m}\;\min_{S \subseteq [m]: q(S) \ge w_i} d_i(S). \]

\begin{lemma}
\label{thm:aps-primal-dual-equiv}
\Cref{defn:aps,defn:aps-dual} are equivalent.
\end{lemma}
\begin{proof}
Let $\pAPS_i$ and $\dAPS_i$ be agent $i$'s AnyPrice shares given by
the primal and dual definitions, respectively.
We will show that for any $z \in \mathbb{R}$, $\pAPS_i \ge z$ iff $\dAPS_i \ge z$.
This would prove that $\pAPS_i = \dAPS_i$.

$\dAPS_i \ge z$ iff the following linear program has a feasible solution:
\[ \min_{x \in \mathbb{R}_{\ge 0}^{\Scal_z}} 0
\textrm{ where } \sum_{S \in \Scal_z} x_S = 1
    \textrm{ and } \left(\forall j \in [m], \sum_{S \in \Scal_z: j \in S} x_S = w_i\right). \]
Its dual is
\[ \max_{p \in \mathbb{R}^m, r \in \mathbb{R}} r - w_ip([m])
\textrm{ where } p(S) \ge r \textrm{ for all } S \in \Scal_z. \]
The dual LP is feasible since $(0, 0)$ is a solution.
Furthermore, if $(p, r)$ is feasible for the dual LP,
then $(\alpha p, \alpha r)$ is also feasible, for any $\alpha \ge 0$.
Hence, by strong duality of LPs, the primal LP is feasible iff
all feasible solutions to the dual have objective value at most 0.

For a given $p$, the optimal $r$ is $\min_{S \in \Scal_z} p(S)$.
Hence, the dual LP is bounded iff for all $p \in \mathbb{R}^m$,
\[ \min_{S \in \Scal_z} p(S) \le w_ip([m]). \]
Furthermore,
\begin{align*}
& \forall p \in \mathbb{R}^m, \min_{S \in \Scal_z} p(S) \le w_ip([m])
\\ &\iff \forall p \in \mathbb{R}^m, \exists S \subseteq [m] \textrm{ such that }
    p(S) \le w_ip([m]) \textrm{ and } v_i(S) \ge z
\\ &\iff \left(\min_{p \in \mathbb{R}^m} \max_{S \subseteq [m]: p(S) \le w_ip([m])} v_i(S)\right) \ge z
\\ &\iff \pAPS_i \ge z.
\end{align*}
Hence, $\dAPS_i \ge z \iff \pAPS_i \ge z$.
\end{proof}
